% !TEX program = pdflatex
\documentclass[11pt,a4paper]{article}

% =========================================================================
% REGISTRO DE INTERESADOS
% Proyecto: CasaScore (AlquilaPuno)
% Gestión de Proyectos TICS (EST315) -- Semestre 2026-II -- Team Sobrevivientes
% =========================================================================

% --- PAQUETES DE CONFIGURACIÓN ---
\usepackage[utf8]{inputenc}
\usepackage[T1]{fontenc}
\usepackage[spanish,es-tabla]{babel}
% Compatibilidad con babel-spanish reciente (flechas de TikZ y signo de porcentaje)
\AtBeginDocument{\shorthandoff{<>}\spanishplainpercent}
\usepackage{geometry}
\geometry{
    top=3.6cm,
    bottom=2.3cm,
    left=2.2cm,
    right=2.2cm,
    % PARTE MODIFICABLE: headheight debe ser mayor que la altura del logo más alto (2.56cm = 73pt).
    headheight=80pt,
    headsep=15pt,
    footskip=20pt
}

\usepackage{graphicx}
\usepackage{fancyhdr}
\usepackage[table]{xcolor}
\usepackage{titlesec}
\usepackage{booktabs}
\usepackage{tabularx}
\usepackage{array}
\usepackage{amsmath}
\usepackage{enumitem}
\usepackage{microtype}
\usepackage{amssymb}
% --- RECURSOS VISUALES: GRÁFICOS Y CÓDIGO DE COLORES ---
\usepackage{tikz}
\usetikzlibrary{positioning,calc}
\definecolor{GrisFila}{RGB}{246, 248, 251}
\definecolor{CuadCerca}{RGB}{150, 30, 30}
\definecolor{CuadSatisf}{RGB}{180, 130, 20}
\definecolor{CuadInfo}{RGB}{10, 45, 90}
\definecolor{CuadMonit}{RGB}{110, 115, 125}
\newcommand{\etiqueta}[2]{{\setlength{\fboxsep}{2pt}\colorbox{#1}{\textcolor{white}{\scriptsize\bfseries #2}}}}
\newcommand{\cerca}{\etiqueta{CuadCerca}{Gestionar de cerca}}
\newcommand{\satisf}{\etiqueta{CuadSatisf}{Mantener satisfecho}}
\newcommand{\informado}{\etiqueta{CuadInfo}{Mantener informado}}
\newcommand{\monit}{\etiqueta{CuadMonit}{Monitorear}}
\usepackage{needspace}

\usepackage{hyperref}

% Logos: funciona a 2 o 3 niveles de profundidad dentro de docs/latex/
\graphicspath{{../../assets/img/}{../../../assets/img/}}

% --- PALETA INSTITUCIONAL FIEI / UNA PUNO ---
\definecolor{AzulUNA}{RGB}{10, 45, 90}
\definecolor{DoradoFIEI}{RGB}{180, 130, 20}
\definecolor{GrisOscuro}{RGB}{50, 50, 50}
\definecolor{AzulClaro}{RGB}{228, 236, 246}

% --- CONFIGURACIÓN DE LISTAS ---
\setlist{topsep=2pt, itemsep=2pt, parsep=1pt}

% --- COLUMNAS DE TABLA ---
\newcolumntype{L}{>{\raggedright\arraybackslash}X}
\newcolumntype{P}[1]{>{\raggedright\arraybackslash}p{#1}}
\newcolumntype{C}[1]{>{\centering\arraybackslash}p{#1}}
\newcolumntype{Q}[1]{>{\raggedleft\arraybackslash}p{#1}}
\newcolumntype{R}{>{\raggedleft\arraybackslash}X}
\renewcommand{\arraystretch}{1.25}

% --- FORMATO DE ENLACES ---
\hypersetup{
    colorlinks=true,
    linkcolor=AzulUNA,
    urlcolor=AzulUNA,
    citecolor=AzulUNA,
    pdftitle={Registro de Interesados -- CasaScore (AlquilaPuno)},
    pdfauthor={Team Sobrevivientes}
}

% --- FORMATO DE TÍTULOS ---
\titleformat{\section}{\color{AzulUNA}\normalfont\large\bfseries}{\thesection.}{0.7em}{}[\color{AzulUNA}\titlerule]
\titleformat{\subsection}{\color{AzulUNA}\normalfont\normalsize\bfseries}{\thesubsection.}{0.5em}{}
\titlespacing*{\section}{0pt}{9pt}{4pt}
\titlespacing*{\subsection}{0pt}{6pt}{2pt}

% --- ENCABEZADO Y PIE DE PÁGINA (fancyhdr) ---
\pagestyle{fancy}
\fancyhf{}

% ===== TAMAÑO DE LOS LOGOS (PARTE MODIFICABLE) =====
% logo_fiei.png  : height=1.45cm (lado izquierdo)
% logo_equipo.png: height=2.56cm (lado derecho; 25 % más grande porque la imagen tiene más borde blanco)
% Para cambiar el tamaño, modifique solo el valor height=... de cada \includegraphics.
% Si el logo más alto supera 2.6cm, aumente también headheight y top en la configuración de página.
\fancyhead[L]{%
    \includegraphics[height=1.45cm,keepaspectratio]{logo_fiei.png}%
}
\fancyhead[C]{%
    \scriptsize \textbf{\textcolor{AzulUNA}{UNIVERSIDAD NACIONAL DEL ALTIPLANO DE PUNO}}\\[2pt]
    \tiny \textcolor{GrisOscuro}{FACULTAD DE INGENIERÍA ESTADÍSTICA E INFORMÁTICA}\\[2pt]
    \scriptsize \textbf{\textcolor{AzulUNA}{Gestión de Proyectos TICS (EST315) \textbar\ Team Sobrevivientes}}\\[2pt]
    \footnotesize \textbf{\textcolor{DoradoFIEI}{REGISTRO DE INTERESADOS}}%
}
\fancyhead[R]{%
    \includegraphics[height=2.56cm,keepaspectratio]{logo_equipo.png}%
}

\renewcommand{\headrulewidth}{0pt}
\renewcommand{\headrule}{}

\fancyfoot[L]{\scriptsize \textcolor{GrisOscuro}{CasaScore (AlquilaPuno)}}
\fancyfoot[C]{\small \thepage}
\fancyfoot[R]{\scriptsize \textcolor{GrisOscuro}{Puno, Perú -- Semestre 2026-II}}
\renewcommand{\footrulewidth}{0.4pt}

\begin{document}

\begin{center}
    {\Large \textbf{\textcolor{AzulUNA}{REGISTRO DE INTERESADOS}}}\\[3pt]
    {\normalsize \textbf{CasaScore (AlquilaPuno)}}\\[4pt]
    \footnotesize \textbf{Fecha:} viernes 18/09/2026 (semana 4) \quad \textbar\quad \textbf{Versión:} 1.0 \quad \textbar\quad \textbf{Fase:} Planificación
\end{center}

\vspace{0.1cm}

\section{Datos Generales}
\noindent\begin{tabularx}{\textwidth}{@{}P{4.2cm} L@{}}
\toprule
\textbf{Proyecto:} & CasaScore (aplicación AlquilaPuno) \\
\textbf{Responsable:} & Clyde Neil Paricahua Pari -- Project Manager \\
\textbf{Proceso PM4DEV:} & Gestión de interesados (Planificar -- Hacer -- Verificar -- Adaptar) \\
\textbf{Insumos:} & Acta de Constitución, especificación de requerimientos y viabilidad del proyecto \\
\bottomrule
\end{tabularx}

\section{Método de Clasificación}
Un interesado es toda persona u organización que puede influir en el proyecto o verse afectada por él. Cada interesado se califica con dos valores de 1 a 5:
\begin{itemize}[leftmargin=*]
    \item \textbf{Poder:} capacidad de cambiar el rumbo del proyecto (aprobar, financiar, usar o dejar de usar el producto).
    \item \textbf{Interés:} cuánto le importan los resultados del proyecto.
\end{itemize}
Un valor de 4 o 5 se considera alto. La combinación de ambos ubica al interesado en uno de cuatro cuadrantes, y cada cuadrante define cómo se le involucra:

\vspace{0.1cm}
{\small
\noindent\begin{tabularx}{\textwidth}{@{}P{3.4cm} P{3.2cm} L@{}}
\toprule
\rowcolor{AzulClaro}
\textbf{Estrategia} & \textbf{Poder / Interés} & \textbf{Qué significa} \\
\midrule
\cerca & Alto / Alto & Participa en las decisiones y en las pruebas; se le consulta en cada sprint. \\
\satisf & Alto / Bajo & Se atienden sus condiciones para que apoye el proyecto, sin cargarlo de detalles. \\
\informado & Bajo / Alto & Recibe avances periódicos y se escuchan sus opiniones. \\
\monit & Bajo / Bajo & Se revisa de vez en cuando por si cambia su situación. \\
\bottomrule
\end{tabularx}
}

\section{Registro de Interesados}
El arrendatario no se limita a estudiantes: es cualquier persona natural que busque cuarto en la ciudad de Puno, residente en el Perú. Por ahora no se atiende a usuarios extranjeros.

\vspace{0.1cm}
{\small
\setlength{\tabcolsep}{4pt}
\rowcolors{2}{white}{GrisFila}
\noindent\begin{tabularx}{\textwidth}{@{}C{0.6cm} P{2.9cm} C{1.25cm} L C{1.0cm} C{1.1cm} C{2.9cm}@{}}
\toprule
\rowcolor{AzulClaro}
\textbf{ID} & \textbf{Interesado} & \textbf{Tipo} & \textbf{Interés principal} & \textbf{Poder} & \textbf{Interés} & \textbf{Estrategia} \\
\midrule
I1 & Team Sobrevivientes (sponsor y dueño) & Interno & Producto viable para concursos y uso real & 5 & 5 & \cerca \\
I2 & Arrendatarios estudiantes (UNA Puno e institutos) & Externo & Cuarto seguro, barato y cerca de su centro de estudios & 4 & 5 & \cerca \\
I3 & Propietarios y arrendadores de Puno & Externo & Publicar fácil y llegar a inquilinos serios & 4 & 4 & \cerca \\
I4 & Arrendatarios trabajadores y familias & Externo & Cuarto cerca de su trabajo o de los colegios & 3 & 4 & \informado \\
I5 & Docente asesor & Externo & Aplicación correcta de la gestión de proyectos & 4 & 3 & \satisf \\
I6 & Organizadores de concursos de la UNA Puno y ferias & Externo & Proyectos innovadores y bien sustentados & 4 & 2 & \satisf \\
I7 & Anunciantes locales & Externo & Llegar a personas que buscan cuarto & 2 & 3 & \monit \\
I8 & Proveedores de servicios (Supabase, Google Play) & Externo & Cumplimiento de sus términos de uso & 3 & 1 & \monit \\
\bottomrule
\end{tabularx}
}

\section{Matriz Poder--Interés}
\begin{center}
\begin{tikzpicture}[x=1.45cm, y=1.05cm, font=\footnotesize]
    \fill[CuadCerca!12]  (3.5,3.5) rectangle (5.5,5.8);
    \fill[CuadSatisf!14] (0.5,3.5) rectangle (3.5,5.8);
    \fill[CuadInfo!10]   (3.5,0.5) rectangle (5.5,3.5);
    \fill[CuadMonit!12]  (0.5,0.5) rectangle (3.5,3.5);
    \draw[gray!60] (0.5,0.5) rectangle (5.5,5.8);
    \draw[gray!60, dashed] (3.5,0.5) -- (3.5,5.8);
    \draw[gray!60, dashed] (0.5,3.5) -- (5.5,3.5);
    \foreach \v in {1,...,5} {
        \node[below, text=GrisOscuro] at (\v,0.5) {\v};
        \node[left, text=GrisOscuro] at (0.5,\v) {\v};
    }
    \node[below=14pt, font=\small\bfseries, text=AzulUNA] at (3,0.5) {Interés};
    \node[rotate=90, font=\small\bfseries, text=AzulUNA] at (-0.1,3) {Poder};
    \node[font=\scriptsize\bfseries, text=CuadCerca]  at (4.5,5.6) {Gestionar de cerca};
    \node[font=\scriptsize\bfseries, text=CuadSatisf] at (2.0,5.6) {Mantener satisfecho};
    \node[font=\scriptsize\bfseries, text=CuadInfo]   at (4.5,0.75) {Mantener informado};
    \node[font=\scriptsize\bfseries, text=CuadMonit]  at (2.0,0.75) {Monitorear};
    \foreach \x/\y/\n/\c in {5/5/I1/CuadCerca, 5/4/I2/CuadCerca, 4/4/I3/CuadCerca, 4/3/I4/CuadInfo, 3/4/I5/CuadSatisf, 2/4/I6/CuadSatisf, 3/2/I7/CuadMonit, 1/3/I8/CuadMonit} {
        \fill[\c] (\x,\y) circle (0.13);
        \node[right=3pt, font=\scriptsize\bfseries] at (\x,\y) {\n};
    }
\end{tikzpicture}
\end{center}

\Needspace{12\baselineskip}
\section{Plan de Involucramiento}
{\small
\rowcolors{2}{white}{GrisFila}
\noindent\begin{tabularx}{\textwidth}{@{}C{0.6cm} L P{2.4cm} P{2.3cm}@{}}
\toprule
\rowcolor{AzulClaro}
\textbf{ID} & \textbf{Acciones} & \textbf{Frecuencia y canal} & \textbf{Responsable} \\
\midrule
I1 & Aprobar actas, planes y cambios mayores; revisar avances, riesgos y gastos. & Martes y viernes, reunión de equipo & Clyde N. Paricahua Pari \\
I2 & Entrevistas para definir pesos y penalizaciones; piloto con 10 estudiantes; encuesta de satisfacción. & Semanas 6, 13 y 15; Google Forms & Viviana M. Quisbert Quispe \\
I3 & Visitas para levantar cuartos; guía corta y acompañamiento para publicar; encuesta sobre el formulario. & Semanas 6 y 13; visita y WhatsApp & Clyde N. Paricahua Pari \\
I4 & Mensajes con los avances y una invitación a probar la app al cierre del piloto. & Semanas 13 y 16; redes sociales & Viviana M. Quisbert Quispe \\
I5 & Presentar cada entregable del curso y aplicar sus observaciones. & Cada entrega del curso & Clyde N. Paricahua Pari \\
I6 & Revisar bases de concursos y preparar la postulación con el MVP. & Semanas 16 y 17 & Ruth K. Apaza Solis \\
I7 & Lista de negocios cercanos a centros de estudio para contactar en el año 1. & Semana 17 & Ruth K. Apaza Solis \\
I8 & Revisar cambios en los términos y límites de la capa gratuita. & Mensual & Eliseo Tarqui Ajahuana \\
\bottomrule
\end{tabularx}
}

\Needspace{10\baselineskip}
\section{Necesidades que se Convierten en Requisitos}
{\small
\rowcolors{2}{white}{GrisFila}
\noindent\begin{tabularx}{\textwidth}{@{}P{3.2cm} L P{2.6cm}@{}}
\toprule
\rowcolor{AzulClaro}
\textbf{Interesado} & \textbf{Necesidad expresada} & \textbf{Requisito} \\
\midrule
Arrendatarios (I2, I4) & Estar cerca de su universidad, trabajo o colegio, no solo de la UNA & RF-03 Punto de referencia \\
Arrendatarios (I2, I4) & Saber por qué un cuarto tiene ese puntaje & RF-07 Explicación del puntaje \\
Arrendatarios (I2, I4) & Conocer las reglas de convivencia antes de visitar & RF-02 y penalizaciones \\
Propietarios (I3) & Publicar rápido desde el celular & RF-02 y RNF-05 \\
Equipo (I1) & Operar sin costo durante el semestre & RNF-07 Costo \\
\bottomrule
\end{tabularx}
}

\section{Seguimiento del Registro}
\begin{itemize}[leftmargin=*]
    \item \textbf{Verificar:} cada martes se revisa si algún interesado cambió de poder, de interés o de postura.
    \item \textbf{Adaptar:} al cierre de cada fase se actualiza este registro y su plan de involucramiento.
    \item \textbf{Registro de cambios:} toda actualización se guarda como una nueva versión en el repositorio.
\end{itemize}

\section{Firmas}
\noindent En señal de conformidad con el presente registro, los integrantes del equipo firman a continuación:

\vspace{1.5cm}
\begin{center}
\begin{tabular}{c c c}
    \parbox{6.5cm}{\centering \rule{6.0cm}{0.6pt}\\[5pt] \textbf{Clyde Neil Paricahua Pari}\\[2pt] {\footnotesize Project Manager}}
    & \hspace{0.8cm} &
    \parbox{6.5cm}{\centering \rule{6.0cm}{0.6pt}\\[5pt] \textbf{Ruth Karina Apaza Solis}\\[2pt] {\footnotesize Analista de Planificación}}
    \\[2.0cm]
    \parbox{6.5cm}{\centering \rule{6.0cm}{0.6pt}\\[5pt] \textbf{Eliseo Tarqui Ajahuana}\\[2pt] {\footnotesize Especialista Técnico / Arquitectura}}
    & \hspace{0.8cm} &
    \parbox{6.5cm}{\centering \rule{6.0cm}{0.6pt}\\[5pt] \textbf{Viviana Mary Quisbert Quispe}\\[2pt] {\footnotesize Gestora de Comunicación, Riesgos y Documentación}}
\end{tabular}
\end{center}
\end{document}
